\clearpage
\thispagestyle{empty}
\pdfbookmark[0]{Über die Anwendung des Magnetometers zur Bestimmung der absoluten Deklination. Ausgabe und Quellen}{absolute-declination-1841-sources}
\begin{center}\Large\textbf{Über die Anwendung des Magnetometers zur Bestimmung der absoluten Deklination}\end{center}

Carl Friedrich Gauss’ Abhandlung erschien in den \textit{Resultaten aus den Beobachtungen des magnetischen Vereins} für 1841 als Artikel I. Die KI-Übertragung der vollständigen Abhandlung und beider historischer Anmerkungen in modernes Deutsch mit KI-erstelltem Satz folgt dem deutschen Text in den \textit{Werken}, Bd.~V (Göttingen, 1867), S.~436--443. Der geometrische Gedankengang, die Beobachtungsdaten und die Schreibweise der Zahlen folgen dieser Ausgabe.

\medskip\noindent
Zur Terminologie des Theodoliten siehe Helmut Fuchs und Tatjana Koukal, \textit{Vermessungskunde für Forstwirte}, S.~91--92, Universität für Bodenkultur. Die Darstellung folgt Gauss’ Instrumentengeometrie und seinen Konventionen.

\medskip\noindent
\href{https://gdz.sub.uni-goettingen.de/id/PPN236006339}{Gesammelte deutsche Ausgabe (1867)}\par
\noindent\href{https://doi.org/10.1007/978-3-642-49319-5_21}{Deutscher Nachdruck (1877)}\par
\noindent\href{https://boku.ac.at/fileadmin/data/H03000/H85000/H85700/skripten/857143/Vermessungskunde-FW.pdf}{Fuchs und Koukal, Vermessungskunde für Forstwirte}
\clearpage
